\section{Entailment-Aware Query Planning}
\label{sec-rewrite}

The interval index answers the instance query $\mathrm{inst}(c)$ of
Section~\ref{sec-pre}, but a \sparql user writes \texttt{?x rdf:type $c$} and
receives the stored triples only: ten of the fourteen LUBM queries return no
rows (Section~\ref{sec-exp}).
\sys therefore offers RDFS subclass entailment as a \emph{plan rewrite}.
With \texttt{\textbackslash set entail on}, a type pattern with a constant
class is answered by the interval index, and the optimizer, the joins, and
every other operator are unchanged.

\stitle{Semantics}.
Let $\mathcal{T}^{+}$ be $\mathcal{T}$ together with $(e,\texttt{rdf:type},c')$ for
every $(e,\texttt{rdf:type},c)\in\mathcal{T}$ and every $c'$ with $c\in\desc(c')$.
A triple pattern $(s,\texttt{rdf:type},c)$ with a constant class has as
\emph{entailed answers} its answers over $\mathcal{T}^{+}$; every other pattern
is evaluated over $\mathcal{T}$.
A pattern whose class is a variable, a class the hierarchy does not contain,
and every other predicate are left as they were.
Only \texttt{rdfs:subClassOf} is read: no \texttt{rdfs:subPropertyOf}, domain,
range, or OWL axiom, and no inference over the properties of a query.

\stitle{The rewrite}.
Algorithm~\ref{alg-rewrite} runs on the planner's tree before the optimizer.
It replaces a qualifying scan by an \emph{interval type scan} and records how
many rows it expects, which is the number of type triples whose class label
lies in an interval of $I(c)$, read from prefix sums over the labels in
$O(|I(c)|)$ time.
The optimizer then orders the joins with that estimate, as it does for any
scan.
This matters: from the stored triples alone, \texttt{?x rdf:type Student}
looks empty, and a plan that starts from it starts from a pattern that in fact
returns 75,547 rows.

\begin{example}
Over \texttt{Dog}, \texttt{Cat} $\sqsubseteq$ \texttt{Mammal} $\sqsubseteq$ \texttt{Animal}
and \texttt{Bird} $\sqsubseteq$ \texttt{Animal}, the query
\texttt{SELECT ?x ?n WHERE \{ ?x a ex:Animal . ?x ex:name ?n \}} returns one row
without entailment (the entity typed \texttt{Animal} itself) and four with it.
\texttt{\textbackslash explain} shows the rewritten plan, in which the optimizer
has joined the interval scan to the name scan by index nested loops:
\begin{small}\begin{verbatim}
Projection [?x, ?n]
  IndexNestedLoopJoin
    IntervalTypeScan {?x rdf:type ex:Animal}
    TripleScan {?x ex:name ?n}
\end{verbatim}\end{small}
\end{example}

\begin{algorithm}[t]
\caption{\textsc{RewriteTypeScans}$(\mathit{plan})$}\label{alg-rewrite}
\small
\ForEach{triple-scan leaf $(s,p,o)$ of $\mathit{plan}$}{
  \If{$p$ is the constant \texttt{rdf:type}, $o$ is a constant class $c$ of the hierarchy $H$}{
    mark the scan as an interval type scan of $c$\;
    $\mathit{rows} \gets \sum_{[a,b]\in I(c)} \big(\mathit{cnt}(b)-\mathit{cnt}(a-1)\big)$ \tcp*{prefix sums of type triples by label}
  }
}
\end{algorithm}

\stitle{Execution}.
An interval type scan is still a scan: it yields triples
$(e,\texttt{rdf:type},c)$ for the instances of $c$, and the operators above
it read them as they read stored ones.
With a free subject it takes one range scan of the label tree per interval of
$I(c)$ and returns the distinct instances in ascending order.
With a bound subject, as when an index nested-loop join probes it with each
outer row, it reads the subject's stored types from the SPO tree and tests
each label against $I(c)$ by binary search.
An entailed triple is not stored, so a scan under a transaction locks the
triple set and not the triple.

\begin{lemma}[Soundness and completeness]\label{lem-rewrite}
For every class $c$ of the hierarchy, the interval type scan of $c$ returns
$(e,\texttt{rdf:type},c)$ if and only if this triple belongs to $\mathcal{T}^{+}$.
\end{lemma}
\iffull
\begin{proof}
The triple is in $\mathcal{T}^{+}$ iff $e$ has a stored type $c''$ with
$c''\in\desc(c)$.
The label tree holds $(\po(c''),e)$ for every stored type of every entity, so by
Lemma~\ref{lem-correct} $e$ is returned by the range scans of $I(c)$ iff
$\po(c'')$ lies in an interval of $I(c)$ for some stored type $c''$, that is,
iff $c''\in\desc(c)$.
A bound subject is tested with the same condition on its stored types.
\end{proof}
\fi

\stitle{Freshness}.
The hierarchy and the label tree are a snapshot, kept on a private file of the
same buffer-pool engine.
Every update, load, and backend switch advances a version number, and the first
entailed query after a change rebuilds the snapshot from the store: 0.06\,s
for LUBM(10) and its 43 classes, a few tens of milliseconds for the Gene
Ontology.
Incremental maintenance of the labels with gaps~\cite{agrawal1989} would avoid
the rebuild and is not implemented.
The snapshot is not part of a transaction, so an entailed query sees the
committed state at the time it built.
